% Requires \usepackage{booktabs}
\begin{table}[htbp]
  \centering
  \caption{Predicted NV shifts for each habit at R = 1.5 nm. The first three quantities are COUPLING-FREE: interior pressure, lattice strain and lattice parameter follow from tau and the elastic tensor alone, so none of the spin-strain uncertainty touches them, and the lattice parameter is directly measurable by size-resolved XRD. Everything below them is not. Delta D is given once per published coupling set rather than as a single number with an error bar, because the two sets differ by a factor of 1.65 on the axial channel for identical strain and that disagreement, not statistics, is the dominant uncertainty. The two E rows say different things and the difference is the substantive result: E evaluated at the VOLUME-AVERAGED strain is identically zero for every symmetry-complete habit, because that average is hydrostatic; the median $|$E(x)$|$ is non-zero because the field resolved in position is not, being concentrated at the particle edges. An ensemble therefore sees inhomogeneous BROADENING rather than a resolved splitting, which is what the FWHM rows quantify against a 3 MHz reference linewidth. The dark-layer rows are an assumption about where NVs are optically active, not a result: halving it from 1.0 to 0.5 nm moves the width by a factor of 2-4, more than most of the physics in this table. All quantities scale as 1/R away from this radius (P, strain and Delta D exactly; $|$E$|$ and the widths only approximately, since the edge-to-interior geometry does not rescale exactly), so a 3 nm particle halves every number here. Epistemic level L1 throughout: the facet tau are L2 but the continuum body, the literature elastic tensor, the literature couplings and the assumed depth distribution are not. \textbullet\ Elastic constants C11/C12/C44 = 1076/125/576 GPa are LITERATURE values, not fitted in this project. Source: Diamond cubic stiffness constants as used by Udvarhelyi, Shkolnikov, Gali, Burkard \& Palyi, Phys. Rev. B 98, 075201 (2018), attributed therein to E. Kaxiras, Atomic and Electronic Structure of Solids (Cambridge University Press, 2003). \textbullet\ udvarhelyi2018\_dft: Udvarhelyi et al., PRB 98, 075201 (2018), Table I \textbullet\ barson2017\_scaled: Udvarhelyi tensor rescaled to Barson et al., Nano Lett. 17, 1496 (2017) axial/transverse magnitudes \textbullet\ Delta D depends only on h41 and h43; E only on h15 and h16. A hydrostatic-pressure ODMR calibration constrains the former and says nothing about the latter, so E carries strictly more uncertainty than Delta D. \textbullet\ Positive interior P is compressive (CLAUDE.md sec 2).}
  \label{tab:t4-predicted-D-and-E}
  \begin{tabular}{lrrrr}
    \toprule
    quantity & cube \{100\} & rh. dodec. \{110\} & octahedron \{111\} & Wulff mixture \\
    \midrule
    facet area fractions & \{100\} 1.00 & \{110\} 1.00 & \{111\} 1.00 & \{100\} 0.16, \{110\} 0.39, \{111\} 0.45 \\
    interior P (GPa, + = compressive) & +2.661 & -4.361 & -0.644 & -1.539 \\
    lattice strain (\%) & -0.201 & +0.329 & +0.049 & +0.116 \\
    lattice parameter a (A) & 3.5658 & 3.5847 & 3.5747 & 3.5771 \\
    Delta D (MHz), udvarhelyi2018\_dft & +21.2 & -34.7 & -5.1 & -12.2 \\
    Delta D (MHz), barson2017\_scaled & +35.0 & -57.3 & -8.5 & -20.2 \\
    E from the volume-averaged strain (MHz) & 0.000 & 0.000 & 0.000 & 0.000 \\
    median $|$E(x)$|$ (MHz) at depth 0.5 nm & 16.07 & 8.89 & 1.98 & 13.27 \\
    median $|$E(x)$|$ (MHz) at depth 1.0 nm & 5.39 & 1.40 & 0.62 & 4.62 \\
    median $|$E(x)$|$ (MHz) at depth 1.5 nm & 1.26 & 0.10 & 0.03 & 0.95 \\
    ODMR FWHM (MHz), dark layer 0.5 nm & 37.8 & 2.8 & 4.6 & 19.2 \\
    ODMR FWHM (MHz), dark layer 1.0 nm & 17.5 & 1.8 & 1.1 & 9.9 \\
    fraction with 2$|$E$|$ $>$ linewidth, dark layer 0.5 nm & 0.97 & 0.82 & 0.36 & 0.95 \\
    fraction with 2$|$E$|$ $>$ linewidth, dark layer 1.0 nm & 0.88 & 0.12 & 0.00 & 0.73 \\
    \bottomrule
  \end{tabular}
\end{table}
